\section{Parallel Repetition}
\label{sec:parrep}

\def\tvanch{\hat{\verifier}^\anch}
\def\tdanch{\hat{\decider}^\anch}
\def\tsanch{\hat{\sampler}^\anch}
\def\vanch{\verifier^\anch}
\def\danch{\decider^\anch}
\def\sanch{\sampler^\anch}

In each of the transformations on verifiers presented so far (introspection,
oracularization, and answer reduction), the \emph{soundness gap} of the
resulting verifier is slightly degraded: while the completeness property (i.e.\
the property of having a PCC strategy with success probability $1$) is
preserved, if the starting game $\verifier_n$ has value at most $1 - \eps$, the
resulting game $\verifier_n'$ has value at most $1 - \delta(\eps, n)$ for some
function $\delta$.
In order to apply the compression procedure recursively we need a way to restore
the soundness gap after a sequence of transformations.
We accomplish this using (a modification of) the technique of \emph{parallel
  repetition}.
  This is a transformation of a two-player game $\game$ into another two-player game $\game^k$
  where the verifier samples $k$ independent pairs $(x_1,y_1),\ldots,(x_k,y_k)$ from the
  question distribution $\mu$ of $\game$, sends the tuple $(x_1,\ldots,x_k)$ to the first player
  to obtain an answer tuple $(a_1,\ldots,a_k)$, and sends the tuple $(y_1,\ldots,y_k)$ to the second
  player to obtain an answer tuple $(b_1,\ldots,b_k)$. The players win if and only if $D(x_i,y_i,a_i,b_i) = 1$ 
  for all $i$.  


Intuitively, if the value of $\game$ is $v < 1$, then one would expect the value
of $\game^k$ to decay exponentially with the number of repetitions $k$.
It is not true in general that the value of $\game^k$ is $v^k$, but exponential
decay bounds on the (tensor product) value of parallel-repeated games are known
for specific classes of
games~\cite{jain2014parallel,dinur2015parallel,bavarian2017hardness}.
In particular, it was shown in~\cite{bavarian2017hardness} that the class of
\emph{anchored games} satisfies exponential-decay parallel repetition, and
furthermore every game can be efficiently transformed into an equivalent
anchored game. \hnote{TODO: update to the SICOMP reference when it gets published.}
Put together, this gives a general soundness amplification procedure called
``anchored parallel repetition,'' which we use in our compression procedure to
reset the soundness gap to a fixed constant.

\hnote{added Dec 29, 2021:} We describe the ``anchoring'' transformation of a game $\game$ in more detail. The anchoring of $\game$ is another game $\game_\perp$ where the verifier behaves as follows: it samples a question pair $(x,y)$ from the question distribution $\mu$ of
$\game$, and then for each player independently changes their question to a special symbol $\perp$ (not part of the original game $\game$) with probability $1/2$ to
obtain new questions $(x',y')$. The first player receives question $x'$ and responds with answer $a$; the second player receives question $y'$ and responds with answer $b$. If neither player receives the question $\perp$, then they win if and only if $D(x',y',a,b) = 1$ where $D$ is the decision predicate of $\game$. Otherwise, the players win if and only if the players receiving the question $\perp$ output a specific, fixed answer (e.g. $0$). 

It is easy to see that $\val^*(\game_\perp) = \frac{3}{4} + \frac{1}{4} \val^*(\game)$. In particular, we have that $\val^*(\game_\perp) = 1$ if and only if $\val^*(\game) = 1$. Bavarian, Vidick and Yuen proved the following bound on the entanglement requirements of $\game_\perp^k$, which is the $k$-fold repetition of $\game_\perp$. 

\begin{theorem}[Parallel repetition of anchored games, Theorem 6.1 of~\cite{bavarian2017hardness}]
\label{thm:bvy}
There exists a universal constant $c > 0$ such that for all two-player games $\game$ and for all positive integers $k$, for all $0 < \eps \leq 1$, for all $p$ satisfying
\[
 p > \frac{4}{\eps}\, \exp \left ( - \frac{c \, \eps^{17} \, k}{s} \right)
\]
where $s$ denotes the bit-length of the players' answers in the game $\game$, we have
\[
	\Ent(\game_\perp^k,p) \geq \Ent(\game,1-\eps)~.
\]
\end{theorem}
We note that a corollary of~\Cref{thm:bvy} is that if $\val^*(\game) < 1 - \eps$, then $\val^*(\game_\perp^k) \leq p$ for $p = \frac{4}{\eps}\, \exp \left ( - \frac{c \, \eps^{17} \, k}{s} \right)$. This is because if $\val^*(\game_\perp^k)$ were larger than $p$, there would be a finite upper bound on $\Ent(\game_\perp^k,p)$, but on the other hand $\Ent(\game,1-\eps)$ is by definition infinite, a contradiction. 

In this section, we present the anchoring and parallel repetition transformations on normal form verifiers; starting with a normal form verifier $\verifier = (\sampler,\decider)$, we first apply the anchoring transformation to obtain a normal form verifier $\vanch = (\sanch,\danch)$, then apply parallel repetition to obtain a normal form verifier $\verifier^\rep = (\sampler^\rep,\decider^\rep)$. The main theorem of the section, \cref{thm:repetition}, establishes the properties of the verifier $\verifier^\rep$.




\begin{remark}
The parallel repetition theorems of~\cite{dinur2015parallel,yuen2016parallel}
are also applicable for the purpose of amplification of the soundness gap, but are not
sufficient for us. This is because the (proofs of the) parallel repetition theorems
of~\cite{dinur2015parallel,yuen2016parallel} only imply that $\Ent(\game^k,v)
\geq \log \Ent(\game,1 - \eps)$, which is not sufficient for the recursive compression application; we need a larger entanglement lower bound 
on $\game^k$. \hnote{moved footnote to remark} The weaker entanglement bound is due to the use of the
  ``quantum correlated sampling lemma'' of~\cite{dinur2015parallel}; it is an open question whether there is
  an improved analysis that obtains a better entanglement lower bound.
\end{remark}


\subsection{The anchoring transformation} 
\label{sec:anchoring}

We present the anchoring transformation on normal form verifiers $\verifier = (\sampler,\decider)$,
which produces another normal form verifier $\vanch = (\sanch,\danch)$.
We first define a typed verifier $\tvanch = (\tsanch,
\tdanch)$, and then detype $\tvanch$ using
Lemma~\ref{lem:detyping-verifiers} to obtain $\vanch$.

Define the type set $\type^\anch = \{\Game, \Anchor\}$ and type graph
$G^\anch$ to be the complete graph over $\type^\anch$ along with self-loops at
each vertex. The Turing machine $\tsanch$ is defined as follows:
\begin{enumerate}
	\item On input $(n,\gamestyle{dimension})$, it returns $\sampler(n,\gamestyle{dimension})$.
	\item On input $(n,w,\gamestyle{marginal},j,z,\tvar)$, it returns $\sampler(n,w,\gamestyle{marginal},j)$ if $\tvar = \Game$, and otherwise returns the binary representation of the zero vector in $\F_2^s$ where $s = \sampler(n,\gamestyle{dimension})$.
	\item On input $(n,w,\gamestyle{linear},j,u,y,\tvar)$, it returns $\sampler(n,w,\gamestyle{linear},j,u,y)$ if $\tvar = \Game$, and otherwise returns the binary representation of the zero vector in $\F_2^s$ where $s = \sampler(n,\gamestyle{dimension})$.
	\item On input $(n,w,\gamestyle{factor},j,u,\tvar)$, it returns $\sampler(n,w,\gamestyle{factor},j,u)$ if $\tvar = \Game$, and otherwise returns the binary representation of the zero vector in $\F_2^s$ where $s = \sampler(n,\gamestyle{dimension})$.
\end{enumerate}
Define the Turing machine $\tdanch$ that, on input
$(n,\tvar_\alice,x_\alice,\tvar_\bob,x_\bob,a_\alice,a_\bob)$, if either
$\tvar_\alice$ or $\tvar_\bob$ is equal to the type $\Anchor$, then the decider
accepts as long as the players receiving the $\Anchor$-type answer with $0$ (a player who receives the $\Game$ type can respond with any answer). 
Otherwise, it accepts only if $\decider(n,x_\alice,x_\bob,a_\alice,a_\bob)$
accepts. Finally, define the Turing machines $(\sanch,\danch)$ to be the result of the detyping transformation $\detype(\tvanch)$ from \Cref{lem:detyping-verifiers} on the pair $\tvanch = (\tsanch,\tdanch)$. 

Note that this transformation is well-defined for \emph{any} pair of Turing machines $(\sampler,\decider)$, even those that don't correspond to normal form verifiers! This transformation can be performed in time that is polynomial in the length of the descriptions of $\sampler$ and $\decider$, and always outputs a pair of Turing machines $(\sanch,\danch)$. However, the following proposition establishes that if $\verifier = (\sampler,\decider)$ is furthermore a normal form verifier, then so is $\vanch = (\sanch,\danch)$ with certain completeness, soundness, and complexity properties.






\begin{proposition}
  \label{prop:anchoring}
  If $\verifier = (\sampler, \decider)$ is a normal form verifier, then 
$\vanch = (\sanch,\danch)$ is a normal form verifier satisfying the following
   for all $n \in \N$.  
\begin{enumerate}
  \item \label{enu:anch-completeness} (\textbf{Completeness}) If there is a
    value-$1$ PCC strategy for $\verifier_n$, then there is a value-$1$ PCC
    strategy for $\vanch_n$.
  \item \label{enu:anch-soundness} (\textbf{Soundness}) For all $\eps > 0$, we have $\Ent(\vanch_n,1 - \eps) \geq \Ent \left (\verifier_n, 1 - (4
        \cdot 16^2) \eps \right)$.
\item \label{enu:anch-complexity} (\textbf{Complexity}) The
    time complexities of the verifier $\vanch$ satisfy
    \begin{align*}
      \TIME_{\sanch}(n) & = \poly(\TIME_\sampler(n))\;,\\
      \TIME_{\danch}(n) & = \poly(\TIME_\decider(n))\;.
    \end{align*}
    Furthermore the number of levels of $\sanch$ is $\ell + 2$ and the dimension is $s(n)+8$, where the number of levels and dimension of sampler $\sampler$ is $\ell$ and $s(n)$ respectively.
  \item \label{enu:anch-efficient-computability} (\textbf{Efficient
      computability}) The descriptions of $\sanch$ and $\danch$
    can be computed in polynomial time from the descriptions of $\sampler$ and
    $\decider$, respectively.
    In particular, the sampler $\sanch$ only depends on the sampler
    $\sampler$.
  \end{enumerate}
\end{proposition}

\begin{proof}
If $\sampler$ is an $\ell$-level sampler with dimension $s(n)$, then by construction $\tsanch$ is a $(\type^\anch, G^\anch)$-typed sampler, with
field size $q(n) = 2$ and the same dimension $s(n)$. It has the following CL functions. Fix an integer $n \in \N$.
Let $V = \F_2^{s(n)}$ denote the ambient space of $\sampler$ on index $n$.
Let $L^\alice, L^\bob: V \to V$ denote the CL functions of $\sampler$ on index
$n$.
For $w\in\{\alice,\bob\}$ the associated CL functions $\{ L^{\anch,\, w}_\tvar
\}$ of $\tsanch$ are
\begin{equation*}
	L^{\anch,\, w}_\tvar =
  \begin{cases}
		 \; L^w & \quad \text{ if } \tvar = \Game \;, \\
		 \; 0 & \quad \text{ if } \tvar = \Anchor \;.
	\end{cases}
\end{equation*}
Intuitively, when the type $\tvar$ sampled for player $w$ is $\Game$, then they
receive a question $L^w(z)$ as they would according to $\sampler$.
Otherwise if $\tvar = \Anchor$, then their question is the zero string.
Thus if $L^w$ is an $\ell$-level CL function, then $L^{\anch,\, w}_\Game$ is
also an $\ell$-level CL function, and $L^{\anch,\, w}_\Anchor$ is a $0$-level CL
function.

  We analyze the completeness, soundness, and complexity properties of the typed
  verifier $\tvanch$; the corresponding properties of the
  detyped verifier $\vanch$ follow from
  \cref{lem:detyping-verifiers} and the fact that the type set
  $\type^\anch$ has size $2$.

  Fix an index $n \in \N$.
  For the completeness property, let $\strategy$ be a value-$1$ PCC strategy for
  $\verifier_n$.
  We define a value-$1$ PCC strategy $\strategy^\anch$ for
  $\tvanch_n$: whenever a player receives the $\Anchor$ type as
  a question type, they perform the trivial measurement (i.e.\ measure the
  identity operator), and output $0$.
  Otherwise, the player performs the same measurement as in $\strategy$.
  This is clearly value-$1$ and PCC.
  \Cref{enu:anch-completeness} follows from this and
  \cref{lem:detyping-verifiers}.

  For the soundness property, we observe that if a strategy $\strategy^\anch$
  has value $1 - \eps$ in $\tvanch_n$, then
  \[ 1 - \eps = \frac{3}{4} + \frac{p}{4} \]
  where $p$ is the value of $\strategy^\anch$ in the game $\verifier_n$, and
  $1/4$ is the probability that neither player receives the question
  type $\Anchor$; this follows from the distribution associated with the typed
  sampler $\tsanch$.
  Solving for $p$, this implies that $\Ent(\tvanch_n,1 - \eps) \geq \Ent(\verifier_n,1 - 4\eps)$.
  The soundness property (\Cref{enu:anch-soundness}) then follows from \cref{lem:detyping-verifiers}.
  


  \Cref{enu:anch-complexity,enu:anch-efficient-computability} are
  straightforward and also follow from \cref{lem:detyping-verifiers}.
\end{proof}

\subsection{The parallel repetition transformation} 
\label{sec:anchored-repetition}
Starting with the anchored verifier $\vanch = (\sanch,\danch)$, we 
then perform the parallel repetition transformation.

Fix integers $\lambda, \tau \in \N$, and let $k(n) = (\lambda n)^{(1 + c')\tau}$ where $c' >0$ is the universal constant such that $\TIME_{\danch}(n) \leq c'(\TIME_\decider(n))^{c'}$ for all $n \geq 1$. Define the Turing machine $\sampler^\rep$ (depending on $\lambda,\tau$) as follows. 
\begin{enumerate}
	\item On input $(n,\gamestyle{dimension})$, it computes $s' = \sanch(n,\gamestyle{dimension})$ and $k = k(n)$, interprets $s'$ as a positive integer and outputs $k \cdot s'$. 
	\item On input $(n,w,\gamestyle{marginal},j,z)$, it computes $s' = \sanch(n,\gamestyle{dimension})$ and $k = k(n)$, parses $z = (z_1,\ldots,z_k) \in (\F_2^{s'})^k$, and then outputs
	\[
		\Big ( \sanch(n,w,\gamestyle{marginal},j,z_1), \ldots, \sanch(n,w,\gamestyle{marginal},j,z_k)\Big)~.
	\]
	\item On input $(n,w,\gamestyle{linear},j,u,y)$, it computes $s' = \sanch(n,\gamestyle{dimension})$ and $k = k(n)$, parses $u = (u_1,\ldots,u_k), y = (y_1,\ldots,y_k) \in (\F_2^{s'})^k$, and then outputs
	\[
		\Big ( \sanch(n,w,\gamestyle{linear},j,u_1,y_1), \ldots, \sanch(n,w,\gamestyle{linear},j,u_k,y_k)\Big)~.
	\]
	\item On input $(n,w,\gamestyle{factor},j,u)$, it computes $s' = \sanch(n,\gamestyle{dimension})$ and $k = k(n)$, parses $u = (u_1,\ldots,u_k) \in (\F_2^{s'})^k$, and then outputs
	\[
		\Big ( \sanch(n,w,\gamestyle{factor},j,u_1), \ldots, \sanch(n,w,\gamestyle{factor},j,u_k)\Big)~.
	\]
\end{enumerate}
\hnote{edited this in response to MdS's email, Oct 2, 2022:} 
Intuitively, the Turing machine $\decider^\rep$ behaves as follows. On input
$(n,{x},{y},{a},{b})$, Parse $x,y,a,b$ as $k(n)$-tuples of questions and answers, respectively, and accept if and only if $\danch(n,x_i,y_i,a_i,b_i)$
accepts for all $i \in \{ 1,\ldots,k(n) \}$. One potential issue is that, depending on how this parsing is implemented, the time complexity of the decider could be unbounded (for example, if $x$ is not properly formatted as a $k(n)$-tuple, then $\decider^\rep$ could take time that grows with the length of $x$, which could be much greater than $\TIME_{\danch}(n)$). To avoid this issue, the decider $\decider^\rep$ checks if any component of the tuples $x,y,a,b$ have length larger than $(\lambda n)^\tau$ bits, and if so then it rejects (this includes the case that $x,y,a,b$ are not properly formatted as $k(n)$-tuples). 
Thus, as long as the time complexity $\TIME_\danch(n)$ is at most $(\lambda n)^{\tau c'}$, then the strings $x,y,a,b$ never need to be longer than $k(n) \cdot (\lambda n)^\tau = (\lambda n)^{(1 + 2c')\tau}$ bits long in order to be accepted by the decider.

Note again that this transformation is well-defined for all Turing machines $(\sanch,\danch)$. The next theorem shows that, if $(\sanch,\danch)$ is the anchoring of a normal form verifier $\verifier = (\sampler,\decider)$, then $\verifier^\rep = (\sampler^\rep,\decider^\rep)$ is also a normal form verifier with certain completeness, soundness, and complexity properties. This is the main result of the section.






\begin{theorem}[Anchored parallel repetition of normal form verifiers]
  \label{thm:repetition}
  There exist universal constants $c,c' > 0$ and a polynomial-time Turing machine
  $\ComputeParrepVerifier$ that given input a tuple $(\verifier,\lambda,\tau)$ where $\verifier = (\sampler,\decider)$ is a pair of Turing machines and $\lambda,\tau$ are positive integers, 
  outputs a pair of Turing machines $\verifier^\rep = (\sampler^\rep,
  \decider^\rep)$ 
  such that the following holds.
	If $\verifier = (\sampler,\decider)$ is an $\ell$-level normal form verifier, then the output $\verifier^\rep$ 
	is a normal form verifier satisfying, for all $n \in \N$, letting $k(n) = (\lambda n)^{(1+c')\tau}$, 
\begin{enumerate}
	\item \label{enu:pr-completeness} (\textbf{Completeness}) If $\verifier_n$ has
    a value-$1$ PCC strategy and $\TIME_\decider(n) \leq (\lambda n)^\tau$, then $\verifier^{\rep}_n$ has a value-$1$ PCC
    strategy.

  \item \label{enu:pr-soundness} (\textbf{Soundness}) For all $\eps > 0$, for
    all
    \begin{equation*}
      p > \frac{4}{\eps} \exp \left ( - c \, \eps^{17} \, k(n)/(\lambda n)^{\tau c'} \right), \end{equation*}
	we have $\Ent(\verifier^\rep_n, p) \geq \Ent(\verifier_n, 1 - \eps)$.


  \item \label{enu:pr-complexity} (\textbf{Complexity}) The repeated verifier
    $\verifier^\rep$ is $(\ell+2)$-level and has time complexities
    \begin{align*}
      \TIME_{\sampler^\rep}(n) & = O(k(n) \cdot \TIME_\sampler(n))\;,\\
      \TIME_{\decider^\rep}(n) & = O \Big(k(n) \cdot \max \Big(\TIME_\decider(n) ,  (\lambda n)^\tau \Big) \Big)\;.
    \end{align*}
  \end{enumerate}
  Furthermore, the repeated sampler $\sampler^\rep$ only depends on $\sampler$ and the parameters $\lambda,\tau$.
\end{theorem}

\begin{proof}

The Turing machine $\ComputeParrepVerifier$, given input $(\verifier,\lambda,\tau)$, first
computes the description of the Turing machines $(\tsanch,\tdanch)$ and then their detyping $(\sanch,\danch)$ 
as described in \Cref{sec:anchoring}. Then, it computes the pair of Turing machines $(\sampler^\rep,\decider^\rep)$ as described above. 
Clearly $\ComputeParrepVerifier$ runs in polynomial time. 


Suppose that $\sampler$ is an $\ell$-level sampler with dimension $s(n)$. Then by \Cref{prop:anchoring}, the number of levels of $\sanch$ is $\ell' = \ell+2$, and its dimension is $s'(n) = s(n)+8$. Then by construction the Turing machine $\sampler^\rep$ is a sampler with the following properties. It has field size $q(n) = 2$ 
and it has dimension $s^\rep(n) = k(n) s'(n)$.
We treat the ambient space $V^\rep$ of $\sampler^\rep$ on index $n$ as the
$k(n)$-fold direct sum of the ambient space $V^\anch$ of $\sanch$ on index $n$.
For all integers $n \in \N$, $w \in \AB$, the CL functions $L^{\rep,\,
  w}: V^\rep \to V^\rep$ of $\sampler^\rep$ are as follows:
\begin{equation*}
	L^{\rep,\, w} = \bigoplus_{i = 1}^{k(n)} L^w\;,
\end{equation*}
where $L^\alice,L^\bob:V^\anch \to V^\anch$ are the CL functions of the sampler $\sampler^\anch$
on index $n$. The CL functions $L^{\rep,\, w}$ are $\ell$-level; the $j$-th
factor spaces $V^{\rep,\, w}_{j,\, u}$ of $L^{\rep,\, w}$ are defined as
\begin{equation*}
	V^{\rep,\, w}_{j,\, u} = \bigoplus_{i = 1}^{k(n)} V^w_{j,\, u_i}
\end{equation*}
for all $u = (u_1,\ldots,u_{k(n)})$ where $V^w_{j,\, u_i}$ is the $j$-th factor
space of $L^w$ with prefix $u_i \in V^\anch$.







The number of levels of $\sampler^\rep$ is the same as $\sanch$, which is $\ell' = \ell + 2$. 

Thus, assuming that $\verifier = (\sampler,\decider)$ is a normal form verifier, we obtain that $\verifier^\rep = (\sampler^\rep,\decider^\rep)$ is a normal form verifier, and thus defines an infinite family of games $(\verifier_n^\rep)_{n \in \N}$. 



	\Cref{enu:pr-completeness} follows from the following straightforward
  observation: if $\strategy = (\psi,A,B)$ is a value-$1$ PCC strategy for
  $\verifier_n$ and $\TIME_\decider(n) \leq (\lambda n)^\tau$, then it must be
  that the question and answer lengths in the strategy $\strategy$ are at most $(\lambda n)^\tau$. Now
  consider the strategy where the players share $k(n)$ copies of
  $\ket{\psi}$, and for the $i$-th instance of the game $\vanch_n$, the
  players use strategy $\strategy$ on the $i$-th copy of $\ket{\psi}$ (and
  performing the identity measurement whenever they receive the $\Anchor$ type). 
  The total length of the questions and answers in the repeated game are at most $k(n) \cdot (\lambda n)^\tau$. 
  It is straightforward to check that this strategy has value $1$ and is PCC.

	To show \Cref{enu:pr-soundness} we apply \Cref{thm:bvy} and \Cref{prop:anchoring}. 
  The exponential decay bound on the value of $\verifier^\rep_n$ presented
  in \Cref{thm:bvy} depends on the answer length of the players in
  the game $\verifier_n^\anch$.
  By \Cref{def:normal-game}, this answer length is at most
  $\TIME_{\decider^\anch}(n) = c' (\TIME_\decider(n))^{c'} \leq (\lambda n)^{\tau c'}$ for some universal constant $c' > 0$ so the claimed bound follows.
	
	\Cref{enu:pr-complexity} follows from the fact that computing the direct sum
  of $k(n)$ CL functions and factor spaces of the ``single-copy'' sampler
  $\sanch$ requires $k(n)$ times the complexity of the ``single-copy''
  sampler, and the complexity of $\sanch$ follows from
  Proposition~\ref{prop:anchoring}.
	The time complexity of the decider follows from the repeated decider having to
  run $k(n)$ instances of the decider $\danch$ and take the logical AND
  of the $k(n)$ decider outputs, as well as the fact that the decider rejects (and thus halts) in the case that the question and answer tuples are improperly formatted. 
  The complexity of $\danch$ follows from
  Proposition~\ref{prop:anchoring} which in turn runs an instance of the
  original decider $\decider$.
  Since the CL functions of the sampler $\sanch$ are $(\ell+2)$-level
  (by Proposition~\ref{prop:anchoring}), and taking the direct sum of CL
  functions does not increase the number of levels (by
  Lemma~\ref{lem:cl-func-prod}), the CL functions $L^{\rep,\, w}$ are
  $(\ell+2)$-level.
  
  The ``Furthermore'' part of the theorem can be easily verified by inspecting the construction of the samplers $\sampler^\rep$ and $\sanch$. 
\end{proof}


 
\section{Gap-Preserving Compression}
